% !TeX program = xelatex
\documentclass[11pt,a4paper]{article}
\usepackage[utf8]{inputenc}
\usepackage[T1]{fontenc}
\usepackage[english,ngerman]{babel}   % Deutsche Rechtschreibung
\usepackage{amsmath,amssymb,amsfonts,mathtools,amsthm}
\usepackage{geometry}
\geometry{left=1.25cm,right=1.25cm,top=1.25cm,bottom=3.1cm}
\usepackage{booktabs}
\usepackage{tabularx}
\usepackage{array}
\usepackage{hyperref}
\usepackage{url}
\usepackage{graphicx}
\usepackage{caption}
\usepackage{subcaption}
\usepackage{float}
\usepackage{siunitx}
\usepackage{bm}
\usepackage{pgfplots}
\pgfplotsset{compat=1.18}
\usepackage{xcolor}
\usepackage{titlesec}
\usepackage{fancyhdr}
\usepackage{microtype}
\usepackage{fontspec}
\usepackage{tcolorbox}
\usepackage{enumitem}

% YUCT standard fonts
\setmainfont{Calibri}
\setsansfont{Segoe UI}[
BoldFont = Segoe UI Bold,
Ligatures = TeX
]

% YUCT color scheme
\definecolor{skyblue}{RGB}{0,102,204}
\definecolor{darkblue}{RGB}{0,0,139}
\definecolor{gold}{RGB}{255,215,0}

% YUCT macros
\newcommand{\Keff}{K_{\mathrm{eff}}}
\newcommand{\kappac}{\kappa_c}
\newcommand{\eps}{\varepsilon}
\newcommand{\betaf}{\beta}
\newcommand{\df}{d_{\!f}}
\newcommand{\Sodd}{S_{\mathrm{odd}}}
\newcommand{\Seven}{S_{\mathrm{even}}}
\newcommand{\Mpl}{M_{\mathrm{Pl}}}
\newcommand{\Lpl}{l_{\mathrm{Pl}}}
\newcommand{\tpl}{t_{\mathrm{Pl}}}
\newcommand{\Lmin}{\ell_{\min}}
\newcommand{\LUV}{\Lambda_{\mathrm{UV}}}

% Theorem environments (German)
\newtheorem{theorem}{Theorem}[section]
\newtheorem{proposition}[theorem]{Proposition}
\newtheorem{definition}[theorem]{Definition}
\theoremstyle{remark}
\newtheorem{remark}[theorem]{Bemerkung}

% Section formatting
\titleformat{\section}{\LARGE\bfseries\color{darkblue}}{\thesection}{1em}{}
\titleformat{\subsection}{\Large\bfseries\color{darkblue}}{\thesubsection}{1em}{}
\titleformat{\subsubsection}{\normalsize\bfseries\color{darkblue}}{\thesubsubsection}{1em}{}

% Headers and footers (angepasst an das russische Template, aber auf Deutsch)
\fancypagestyle{firstpage}{%
	\fancyhf{}%
	\renewcommand{\headrulewidth}{-5pt}%
	\renewcommand{\footrulewidth}{-5pt}%
	\fancyfoot[C]{%
		\begin{minipage}[t]{\textwidth}
			\centering
			\textcolor{gold}{\rule{\textwidth}{1.6pt}}\\[5pt]
			{\small \textsuperscript{1}Yakushev Research, \url{https://yuct.org/}} \hfill
			{\small YPSDC-Protokoll: \url{https://ypsdc.com/}}\\[-2pt]
			\textcolor{gold}{\rule{\textwidth}{1.6pt}}\\[5pt]
			\textbf{EINHEITLICHE KOORDINATIONSTHEORIE JAKUSCHEW 2026} \hfill \thepage\\[10pt]
		\end{minipage}%
	}%
}

\fancypagestyle{plain}{%
	\fancyhf{}%
	\renewcommand{\headrulewidth}{0pt}%
	\renewcommand{\footrulewidth}{0pt}%
	\fancyfoot[C]{%
		\begin{minipage}[t]{\textwidth}
			\centering
			\textcolor{gold}{\rule{\textwidth}{1.6pt}}\\[4pt]
			\textbf{EINHEITLICHE KOORDINATIONSTHEORIE JAKUSCHEW 2026} \hfill \thepage\\[4pt]
		\end{minipage}%
	}%
}

\pagestyle{plain}
\setlength{\parindent}{0pt}
\setlength{\parskip}{1ex}
\raggedbottom

\hypersetup{
	colorlinks=true,
	linkcolor=blue,
	citecolor=blue,
	urlcolor=blue,
}

% Custom header (angepasst: rechte Seite zeigt jetzt "Anhang UVD" vor dem Doppelpunkt)
\newcommand{\makeheader}{%
	\noindent
	\begin{minipage}{0.17\textwidth}
		\raggedright
		{\fontspec{Segoe UI}[BoldFont=Segoe UI Bold]\bfseries\color{skyblue}\fontsize{46}{46}\selectfont YUCT}%
	\end{minipage}%
	\hspace{30pt}
	\begin{minipage}{0.32\textwidth}
		\raggedright
		{\sffamily\fontsize{11}{13}\selectfont JAKUSCHEW\\ EINHEITLICHE\\ KOORDINATIONSTHEORIE}%
	\end{minipage}%
	\hfill
	\begin{minipage}{0.38\textwidth}
		\raggedleft
		{\sffamily\bfseries Anhang UVD}\\ 
		{\sffamily Version 1.0 / Juni 2026}\\ 
		{\sffamily\footnotesize \url{https://doi.org/10.5281/zenodo.18444598}}%
	\end{minipage}
	\par\vspace{5pt}
	\noindent\textcolor{gold}{\rule{\textwidth}{1.6pt}}
	\par\vspace{0pt}
}

\begin{document}
	
	% Title page
	\thispagestyle{firstpage}
	\makeheader
	
	\vspace{0.5cm}
	{\LARGE\bfseries Anhang UVD: Auflösung ultravioletter Divergenzen in YUCT \\
		\Large Endliche Quantenfeldtheorie aus dem Koordinationsformfaktor und der physikalische Skala des Koordinationsnetzwerks \par}
	\vspace{0.3cm}
	
	{\large Alexey V. Yakushev\textsuperscript{1}} \\
	\vspace{0.2cm}
	
	{\color{darkblue}\textbf{\LARGE Zusammenfassung}} \par
	{\large
		\noindent
		Die einheitliche Koordinationstheorie von Jakuschew (YUCT) schlägt vor, dass ultraviolette Divergenzen in der Quantenfeldtheorie durch einen Koordinationsformfaktor \(F(q^2) = \exp[-(q^2/\LUV^2)^{2/3}]\) beseitigt werden, der aus der endlichen Informationskapazität des Vakuums entsteht. In der ursprünglichen Formulierung wurde die UV-Skala \(\LUV\) mit der Planck-Masse identifiziert, \(\LUV = 3\pi\Mpl \approx 1,15\times 10^{20}\) GeV, was zu einer enormen endlichen Strahlungskorrektur der Higgs-Masse führte. Wir korrigieren einen Rechenfehler in der Konstante \(C_0\), die in der skalaren Massenkorrektur auftritt: Der korrekte Wert ist \(C_0 = 3\sqrt{\pi}/(8\sqrt{2}) \approx 0,46999\) und ersetzt den zuvor berichteten Wert \(C_0 \approx 0,626\). Mit dieser Korrektur beträgt der radiative Beitrag zur Higgs-Masse von einem Netzwerk im Planck-Maßstab \(\delta m \approx 1,57\times 10^{18}\) GeV, was ein schwerwiegendes Feinabstimmungsproblem wieder einführt, es sei denn, die bloße Masse wird exakt kompensiert. Wir untersuchen zwei logische Alternativen: (i) Die YUCT-Massenformel \(m = \Mpl(2/3)^L\xi_f\) bestimmt bereits die physikalische Masse, und die Strahlungskorrektur wird durch einen endlichen Gegenterm absorbiert – eine technisch konsistente, aber ästhetisch unbefriedigende Lösung; (ii) Der fundamentale Lagrange-Operator auf der Planck-Skala enthält keinen bloßen Massenterm (Skaleninvarianz), und die gesamte Higgs-Masse wird durch die Strahlungskorrektur erzeugt. In diesem zweiten Szenario verlangt die Selbstkonsistenz, dass die UV-Skala \(\LUV \approx 8,96\) TeV beträgt – ein Wert, der am LHC zugänglich ist. Der Koordinationsformfaktor sagt dann beobachtbare Abweichungen in hochenergetischen Streuquerschnitten im TeV-Maßstab voraus. Wir präsentieren beide Szenarien, ihre mathematische Rechtfertigung und ihre experimentellen Konsequenzen, um YUCT strengen empirischen Tests zu unterziehen.
	}
	
	\vspace{0.2cm}
	\noindent
	{\color{darkblue}\textbf{Schlüsselwörter:}} YUCT, ultraviolette Divergenzen, Koordinationsformfaktor, Hierarchieproblem, Higgs-Masse, Skaleninvarianz, nichtlokale Quantenfeldtheorie, LHC-Phänomenologie.
	\vspace{0.1cm}
	\noindent
	{\color{darkblue}\textbf{Zitierung:}} Yakushev AV. Anhang UVD: Auflösung ultravioletter Divergenzen in YUCT. 2026. DOI: \url{https://doi.org/10.5281/zenodo.18444598} (dieser Band).
	
	\vspace{0.5cm}
	
	% ============================================================
	\section{Einführung}
	\label{sec:intro}
	
	Die einheitliche Koordinationstheorie von Jakuschew (YUCT) bietet eine elegante Auflösung der ultravioletten Divergenzen der Quantenfeldtheorie (QFT). Anstatt einen willkürlichen Impuls-Cut-off oder eine unendliche Renormierung einzuführen, postuliert YUCT, dass das physikalische Vakuum eine endliche „Bandbreite“ für die Übertragung von Quanteninformationen besitzt. Diese Bandbreite wird durch einen \textbf{Koordinationsformfaktor} \(F(q^2)\) quantifiziert, der virtuelle Impulse oberhalb einer charakteristischen Skala \(\LUV\) exponentiell unterdrückt. Die mathematische Struktur dieses Formfaktors wird durch das universelle Fehlergesetz von YUCT bestimmt: \(\varepsilon \propto K_{\mathrm{eff}}^{-\beta}\) mit \(\beta = 2/3\).
	
	In früheren Versionen dieses Anhangs identifizierten wir \(\LUV\) mit der Planck-Masse, \(\LUV = 3\pi\Mpl \approx 1,15\times 10^{20}\) GeV – eine Skala, die aus der minimalen Länge \(\Lmin = \frac{2}{3}\Lpl\) abgeleitet wurde, die vom primären algebraischen Zyklus von YUCT diktiert wird. Mit dieser Identifikation werden alle Schleifenintegrale endlich: Logarithmische Divergenzen werden durch \(\ln(\LUV/m)\) ersetzt, und Potenzdivergenzen verschwinden identisch.
	
	Eine sorgfältige Neuberechnung der skalaren Massenkorrektur offenbart jedoch einen Rechenfehler in der Konstante \(C_0\), die die Größe des endlichen Rests bestimmt. Der korrekte Wert von \(C_0\) ist kleiner als zuvor berichtet, was die Strahlungskorrektur verringert, sie aber immer noch um viele Größenordnungen über der beobachteten Higgs-Masse belässt. Dies zwingt zu einer kritischen Neubewertung der physikalischen Skala \(\LUV\).
	
	Dieser Anhang hat zwei Ziele. Erstens legen wir die korrigierten mathematischen Ableitungen detailliert dar, um sicherzustellen, dass alle Ausdrücke genau und reproduzierbar sind. Zweitens untersuchen wir die physikalischen Implikationen der korrigierten Formeln, einschließlich eines radikalen Vorschlags: Die fundamentale Skala des Koordinationsnetzwerks könnte nicht die Planck-Skala sein, sondern eine viel niedrigere, möglicherweise zugängliche in aktuellen Collider-Experimenten.
	
	% ============================================================
	\section{Der Koordinationsformfaktor und die UV-Skala}
	\label{sec:formfactor}
	
	Der YUCT-Koordinationsformfaktor lautet
	\begin{equation}
		F(q^2) = \exp\!\Bigl[-\Bigl(\frac{q^2}{\LUV^2}\Bigr)^{\!\beta}\,\Bigr], \qquad \beta = \frac{2}{3}.
		\label{eq:F}
	\end{equation}
	Alle Ausbreiter des Standardmodells werden mit \(F(p^2/\LUV^2)\) multipliziert. Die Skala \(\LUV\) wird ursprünglich aus der minimalen Länge \(\Lmin = \frac{2}{3}\Lpl\) über die Compton-Beziehung \(\LUV = 2\pi\hbar c/\Lmin = 3\pi\Mpl\) abgeleitet.
	
	Dieser Formfaktor macht alle Schleifenintegrale im Ultraviolett absolut konvergent. Der Beweis ist einfach: Für jedes Polynom \(P(k)\) in den Schleifenimpulsen fällt das Produkt \(P(k) \cdot \prod_i F(k_i^2/\LUV^2)\) für \(|k| \to \infty\) schneller als jede inverse Potenz von \(|k|\) ab, weil die Exponentialfunktion dominiert.
	
	% ============================================================
	\section{Elektronenselbstenergie: Auflösung einer logarithmischen Divergenz}
	\label{sec:electron}
	
	Die Ein-Schleifen-Selbstenergie des Elektrons in der QED ist in der konventionellen QFT logarithmisch divergent. Mit dem YUCT-Formfaktor wird das euklidische Integral zu
	\begin{equation}
		\Sigma_{\mathrm{YUCT}} \propto \int_0^{\infty} \frac{dk_E}{k_E}\; \exp\!\Bigl[-2\Bigl(\frac{k_E^2}{\LUV^2}\Bigr)^{\!2/3}\,\Bigr].
		\label{eq:elec_int}
	\end{equation}
	Die Substitution \(t = (k_E^2/\LUV^2)^{2/3}\) liefert \(dk_E/k_E = \frac{3}{4} dt/t\), und das Integral ergibt
	\begin{equation}
		I_{\mathrm{log}} = \frac{3}{4} \int_{t_{\min}}^{\infty} \frac{dt}{t}\, e^{-2t} \approx \ln\frac{\LUV}{m_e},
		\label{eq:elec_result}
	\end{equation}
	wobei \(t_{\min} \sim (m_e^2/\LUV^2)^{2/3}\). Die logarithmische Divergenz wird durch den endlichen Logarithmus \(\ln(\LUV/m_e) \approx 53,8\) ersetzt. Die Berechnung ist elementar und benötigt keine Korrektur.
	
	% ============================================================
	\section{Skalare Massenkorrektur: Korrektur der Konstante \(C_0\)}
	\label{sec:scalar}
	
	\subsection{Das Integral und seine Auswertung}
	
	Die Ein-Schleifen-Massenkorrektur für ein Skalarfeld in der \(\phi^4\)-Theorie lautet nach Wick-Rotation und Einsetzen des Formfaktors
	\begin{equation}
		\delta m^2_{\mathrm{YUCT}} = \frac{\lambda}{16\pi^2} \int_0^{\infty} \frac{k_E^3\, dk_E}{k_E^2 + m^2}\; F^2\!\Bigl(\frac{k_E^2}{\LUV^2}\Bigr), \qquad F^2(x) = e^{-2x^{2/3}}.
		\label{eq:scalar_YUCT}
	\end{equation}
	
	Mit der dimensionslosen Variablen \(t = k_E^2/\LUV^2\) und der Entwicklung für \(m \ll \LUV\) erhält man den führenden Term
	\begin{equation}
		\delta m^2_{\mathrm{YUCT}} = \frac{\lambda}{32\pi^2}\,\LUV^2 \int_0^{\infty} e^{-2t^{2/3}} dt + \mathcal{O}(m^2/\LUV^2).
		\label{eq:scalar_lead}
	\end{equation}
	
	Früher wurde berichtet, dass das Integral \(C_0 \equiv \int_0^{\infty} e^{-2t^{2/3}} dt \approx 0,626\) (Version 3.1) sei. Dieser Wert ist falsch. Wir berechnen ihn nun exakt.
	
	\subsection{Exakte Berechnung von \(C_0\)}
	
	\begin{theorem}[Exakter Wert von \(C_0\)]
		\label{thm:C0}
		\[
		C_0 = \int_0^{\infty} e^{-2t^{2/3}} dt = \frac{3\sqrt{\pi}}{8\sqrt{2}} \approx 0,46999.
		\]
	\end{theorem}
	
	\begin{proof}
		Substituiere \(u = 2t^{2/3}\). Dann ist \(t = (u/2)^{3/2}\) und
		\[
		dt = \frac{3}{2} (u/2)^{1/2} \cdot \frac{1}{2} du = \frac{3}{4\sqrt{2}} u^{1/2} du.
		\]
		Daher
		\[
		C_0 = \int_0^{\infty} e^{-u} \cdot \frac{3}{4\sqrt{2}} u^{1/2} du = \frac{3}{4\sqrt{2}} \int_0^{\infty} u^{1/2} e^{-u} du = \frac{3}{4\sqrt{2}} \, \Gamma(3/2).
		\]
		Mit \(\Gamma(3/2) = \frac{\sqrt{\pi}}{2}\) folgt
		\[
		C_0 = \frac{3}{4\sqrt{2}} \cdot \frac{\sqrt{\pi}}{2} = \frac{3\sqrt{\pi}}{8\sqrt{2}} \approx 0,46999.
		\]
	\end{proof}
	
	\subsection{Korrigierte skalare Massenkorrektur}
	
	Mit dem korrekten Wert von \(C_0\) lautet der führende Term der skalaren Massenkorrektur
	\begin{equation}
		\boxed{\delta m^2_{\mathrm{YUCT}} = \frac{\lambda}{32\pi^2}\,\LUV^2 \cdot \frac{3\sqrt{\pi}}{8\sqrt{2}} + \mathcal{O}(m^2/\LUV^2)}.
		\label{eq:scalar_corrected}
	\end{equation}
	
	Für das Higgs-Boson gilt \(\lambda = m_H^2/(2v^2) \approx 0,13\) (mit \(v \approx 246\) GeV). Unter Verwendung der Planck-Skalen-Identifikation \(\LUV = 3\pi\Mpl \approx 1,15\times 10^{20}\) GeV:
	\begin{align}
		\delta m^2_H &\approx \frac{0,13}{32\pi^2} \cdot (1,15\times 10^{20})^2 \cdot 0,46999 \notag\\
		&\approx 2,47 \times 10^{36}\ \mathrm{GeV}^2, \notag\\
		\delta m_H &\approx 1,57 \times 10^{18}\ \mathrm{GeV}.
		\label{eq:planck_correction}
	\end{align}
	
	Dieser Wert übersteigt die beobachtete Higgs-Masse (\(m_H^{\mathrm{exp}} \approx 125\) GeV) um den Faktor \(1,26 \times 10^{16}\). Wenn die bloße Masse im Lagrange-Operator Null ist, wird die physikalische Masse vollständig durch diese Strahlungskorrektur bestimmt, und der vorhergesagte Wert ist offensichtlich unverträglich mit der Beobachtung.
	
	% ============================================================
	\section{Physikalische Implikationen: Zwei Szenarien}
	\label{sec:scenarios}
	
	Die korrigierte Berechnung zwingt uns, das Hierarchieproblem direkt anzugehen. Wir betrachten zwei logisch konsistente Alternativen.
	
	\subsection{Szenario A: Planck-Skalen-Netzwerk mit endlichen Gegentermen}
	
	In diesem Szenario behält YUCT die Identifikation \(\LUV = 3\pi\Mpl\) bei. Die physikalische Masse jedes Teilchens wird durch die YUCT-Massenformel
	\begin{equation}
		m_{\mathrm{phys}} = \Mpl \left(\frac{2}{3}\right)^L \xi_f,
		\label{eq:mass_formula}
	\end{equation}
	gegeben, wobei \(L\) die Koordinationstiefe und \(\xi_f\) ein Resonanzfaktor ist (siehe Anhang~J). Die bloße Masse im Lagrange-Operator ist ein freier Parameter, der so gewählt wird, dass er die Strahlungskorrektur aufhebt und die beobachtete Masse reproduziert.
	
	Mathematisch ist dies völlig konsistent: Die Theorie ist endlich, der erforderliche Gegenterm ist ebenfalls endlich, und es treten niemals Unendlichkeiten auf. Die YUCT-Massenformel, Gl.~(\ref{eq:mass_formula}), liefert die Renormierungsbedingung, die den Gegenterm festlegt.
	
	Physikalisch erfordert dieses Szenario jedoch, dass sich die bloße Masse und die Strahlungskorrektur auf einen Teil in \(10^{16}\) kompensieren – eine Feinabstimmung genau der Art, die das Hierarchieproblem vermeiden soll. YUCT erklärt nicht, warum die bloße Masse so gewählt werden sollte, um diese Kompensation zu erreichen; es ersetzt lediglich die übliche unendliche Subtraktion durch eine enorme endliche.
	
	\subsection{Szenario B: Skaleninvarianter Ursprung und ein Netzwerk im TeV-Maßstab}
	
	Eine alternative, ästhetisch ansprechendere Lösung besteht darin, die Identifikation von \(\LUV\) mit der Planck-Skala aufzugeben. Wenn die fundamentale Theorie auf der Planck-Skala keine bloßen Massenterme enthält (d.~h. skaleninvariant ist), dann werden alle Massen radiativ durch den Koordinationsformfaktor erzeugt. In diesem Fall ist die physikalische Higgs-Masse \emph{gleich} der Strahlungskorrektur:
	\begin{equation}
		m_H^2 = \delta m^2_{\mathrm{YUCT}} = \frac{\lambda}{32\pi^2}\,\LUV^2 \cdot C_0.
		\label{eq:scale_inv}
	\end{equation}
	
	Auflösen nach \(\LUV\) ergibt
	\begin{equation}
		\LUV = \sqrt{\frac{32\pi^2\, m_H^2}{\lambda\, C_0}}.
		\label{eq:LUV_solve}
	\end{equation}
	
	Setzt man das korrigierte \(C_0 = 3\sqrt{\pi}/(8\sqrt{2}) \approx 0,46999\), \(\lambda = 0,13\) und \(m_H = 125\) GeV ein, erhält man
	\begin{equation}
		\boxed{\LUV \approx 8,96\ \mathrm{TeV}}.
		\label{eq:LUV_TeV}
	\end{equation}
	
	Dieser Wert ist eine Vorhersage der Theorie, kein freier Parameter. Er platziert die fundamentale Koordinationsskala direkt in den TeV-Bereich und macht sie experimentell zugänglich.
	
	\subsection{Physikalische Interpretation von Szenario B}
	
	Wenn \(\LUV \approx 9\) TeV, „kristallisiert“ das Koordinationsnetzwerk auf der elektroschwachen Skala, nicht auf der Planck-Skala. Der Formfaktor \(F(q^2) = \exp[-(q^2/\LUV^2)^{2/3}]\) beginnt, virtuelle Impulse oberhalb einiger TeV zu unterdrücken. Dies hat unmittelbare beobachtbare Konsequenzen:
	
	\begin{enumerate}[label=(\arabic*)]
		\item \textbf{Drell-Yan- und Dijet-Querschnitte} am LHC sollten einen Defizit an Ereignissen mit hohen invarianten Massen im Vergleich zu den Vorhersagen des Standardmodells zeigen. Die Unterdrückung wird für \(m_{jj} \gtrsim 5\) TeV merkbar.
		\item \textbf{Paarproduktion von Higgs-Bosonen} und Vektorbosonstreuung sollten eine anomalie Energieabhängigkeit aufweisen, da der Formfaktor die Higgs-Selbstkopplung modifiziert.
		\item \textbf{Abweichungen in elektroschwachen Präzisionsobservablen} sind möglicherweise zu erwarten, aber wahrscheinlich klein, da die LEP-Messungen Skalen weit unterhalb von \(\LUV\) sondieren.
	\end{enumerate}
	
	Entscheidend ist, dass die minimale Länge in diesem Szenario \(\Lmin = 2\pi\hbar c/\LUV \approx 1,38\times 10^{-19}\) m beträgt, was viel größer ist als die Planck-Länge (\(\sim 10^{-35}\) m). Dies impliziert, dass die Raumzeit-Körnung kein quantengravitativer Effekt ist, sondern ein Phänomen der elektroschwachen Skala, möglicherweise verbunden mit dem Higgs-Mechanismus selbst.
	
	\subsection{Die fraktale Brücke zwischen der elektroschwachen Skala und der kosmologischen Konstanten}
	\label{subsec:bridge}
	
	Wenn Szenario B korrekt ist und die fundamentale Koordinationsskala \(\LUV \approx 8,96\) TeV beträgt, ergibt sich eine tiefgreifende Verbindung zur Kosmologie. In YUCT ist das Vakuum kein einheitliches kontinuierliches Medium, sondern eine fraktale Hierarchie von Koordinationsschichten. Die Skala \(\LUV\) charakterisiert die Schicht, auf der die Felder des Standardmodells den Koordinationsformfaktor erfahren. Es existieren jedoch tiefere (infrarote) Schichten, die sukzessive niedrigeren Energieskalen entsprechen. Die kosmologische Konstante, d.~h. die beobachtete Dichte der Dunklen Energie, wird durch die tiefste infrarote Schicht erzeugt, deren Skala \(\Lambda_{\mathrm{IR}}\) durch die Anforderung festgelegt ist, dass die Nullpunktenergie aller Quantenfelder, integriert über die gesamte fraktale Hierarchie, den beobachteten Wert \(\rho_{\Lambda} \approx 2,51 \times 10^{-47}\) GeV\(^4\) reproduziert (siehe Anhang~M).
	
	Die Beziehung zwischen der UV-Skala \(\LUV\) und der IR-Skala \(\Lambda_{\mathrm{IR}}\) wird durch das fundamentale Skalierungsquantum von YUCT bestimmt:
	\begin{equation}
		q = \left(\frac{3}{2}\right)^{1/3} \approx 1,1447.
	\end{equation}
	Die beiden Skalen sind durch eine ganze Zahl \(N\) von Koordinationsschichten getrennt:
	\begin{equation}
		\LUV = \Lambda_{\mathrm{IR}} \cdot q^{N}.
		\label{eq:bridge}
	\end{equation}
	
	Aus der beobachteten kosmologischen Konstanten extrahiert man \(\Lambda_{\mathrm{IR}} \approx 0,0032\) eV (siehe Anhang~M, Gl.~(M.47)). Mit \(\LUV \approx 8,96 \times 10^{12}\) eV erhält man
	\begin{equation}
		N = \frac{\ln(\LUV / \Lambda_{\mathrm{IR}})}{\ln q}
		= \frac{\ln(8,96 \times 10^{12} / 0,0032)}{\frac{1}{3}\ln(3/2)}
		\approx \frac{35,564}{0,135155} \approx 263,1.
		\label{eq:N_layers}
	\end{equation}
	
	Aufgerundet auf die nächste halbe ganze Zahl (wie vom algebraischen Zyklus von YUCT für fermionische Schichten gefordert) erhält man die fundamentale Strukturzahl
	\begin{equation}
		\boxed{N = 263,0}.
	\end{equation}
	
	Dies ist ein bemerkenswertes Ergebnis. Die elektroschwache Skala (die Masse des Higgs-Bosons) und die kosmologische Konstante sind durch genau \(263\) diskrete Schritte der Koordinationshierarchie verbunden, ohne freie Parameter. Die beobachtete Kleinheit der kosmologischen Konstanten ist kein Rätsel; sie spiegelt einfach die enorme Tiefe des Koordinationsnetzwerks wider, das die Teilchenphysik von der Kosmologie trennt.
	
	\textbf{Interpretation.} Jede Schicht der Hierarchie wirkt als „Skalentransformator“: Die Nullpunktenergie einer gegebenen Schicht wird durch den Koordinationsformfaktor der nächsten, tieferen Schicht unterdrückt. Die gesamte Vakuumenergie ist die Summe der Beiträge aller Schichten, wobei jede Schicht um den Faktor \(\beta = 2/3\) kleiner ist als die vorherige. Die geometrische Reihe konvergiert gegen einen endlichen Wert, der genau die beobachtete Dichte der Dunklen Energie ist (siehe Anhang~M). Die ganze Zahl \(N \approx 263\) ist die Anzahl der Schichten, die zwischen der elektroschwachen Skala und der Neutrinomassenskala (der tiefsten physikalisch zugänglichen Schicht) liegen. Dass diese Zahl eine halbe ganze Zahl ist, ist eine nicht-triviale Konsistenzprüfung des YUCT-Rahmens, die Szenario B weiter stützt.
	
	\subsection{Vergleich der beiden Szenarien}
	
	\begin{table}[H]
		\centering
		\caption{Vergleich von Szenario A und Szenario B.}
		\label{tab:scenarios}
		\small
		\begin{tabular}{@{}lcc@{}}
			\toprule
			\textbf{Eigenschaft} & \textbf{Szenario A} & \textbf{Szenario B} \\
			\midrule
			UV-Skala \(\LUV\) & \(3\pi\Mpl \approx 10^{20}\) GeV & \(\approx 8,96\) TeV \\
			Minimale Länge \(\Lmin\) & \(\frac{2}{3}\Lpl \approx 10^{-35}\) m & \(\approx 10^{-19}\) m \\
			Bloße Masse im Lagrange-Operator & Nicht Null (Feinabstimmung) & Null (Skaleninvarianz) \\
			Strahlungskorrektur \(\delta m_H\) & \(10^{18}\) GeV & 125 GeV \\
			Experimentelle Zugänglichkeit & Planck-Skala (unzugänglich) & LHC-Energien (zugänglich) \\
			Feinabstimmung erforderlich & Ja (1 Teil in \(10^{16}\)) & Nein \\
			\bottomrule
		\end{tabular}
	\end{table}
	
	% ============================================================
	\section{Falsifizierbare Vorhersagen von Szenario B}
	\label{sec:predictions}
	
	Wenn Szenario B korrekt ist, können die folgenden quantitativen Vorhersagen am LHC und an zukünftigen Collidern getestet werden:
	
	\begin{enumerate}[label=(\arabic*)]
		\item \textbf{Unterdrückung des Drell-Yan-Querschnitts.} Das Verhältnis des gemessenen zum Standardmodell-Drell-Yan-Querschnitt als Funktion der invarianten Dileptonmasse \(m_{\ell\ell}\) sollte folgen
		\[
		R(m_{\ell\ell}) \approx \exp\!\Bigl[-2\Bigl(\frac{m_{\ell\ell}^2}{\LUV^2}\Bigr)^{\!2/3}\,\Bigr],
		\]
		mit \(\LUV \approx 8,96\) TeV. Eine 10\%-Unterdrückung wird bei \(m_{\ell\ell} \approx 5,4\) TeV vorhergesagt, die bei \(m_{\ell\ell} \approx 10\) TeV auf einen Faktor von \(\sim 2\) anwächst.
		
		\item \textbf{Paarproduktion von Higgs-Bosonen.} Der Formfaktor modifiziert die Higgs-Selbstkopplung bei hohem Impulsübertrag. Die Verteilung der invarianten Di-Higgs-Masse sollte für \(m_{HH} \gtrsim 3\) TeV vom Standardmodell abweichen.
		
		\item \textbf{Fermionenmassen.} Wenn alle Fermionenmassen ebenfalls radiativ erzeugt werden, müssen ihre Yukawa-Kopplungen so gewählt werden, dass die Strahlungskorrekturen das beobachtete Massenspektrum reproduzieren. Dies legt eine Selbstkonsistenzbedingung fest, die die Koordinationstiefen \(L_f\) und die Resonanzfaktoren \(\xi_f\) mit den Strahlungskorrekturen verknüpft. Eine detaillierte Analyse wird auf zukünftige Arbeiten verschoben.
	\end{enumerate}
	
	\subsection{Eine Bemerkung zur Feinstrukturkonstante}
	
	Die Herleitung von \(\alpha^{-1} \approx 137,036\) aus der topologischen Invarianten \(N_{\mathrm{gauge}} = 12\) (Anhang~G) ist unabhängig vom Wert von \(\LUV\). Die Eichkopplungen laufen logarithmisch; ihre Anfangswerte bei \(\LUV\) werden durch die Topologie von \(F_{18}\) festgelegt. In Szenario A ist \(\LUV \sim \Mpl\), und der logarithmische Lauf von der Planck-Skala bis zur elektroschwachen Skala ergibt \(\alpha^{-1} \approx 137,036\). In Szenario B ist \(\LUV \sim 9\) TeV, und der Lauf erfolgt über ein viel kürzeres Intervall. Die topologische Vorhersage für \(\alpha^{-1}\) würde dann auf der TeV-Skala gelten, und der niederenergetische Wert würde durch eine kleine logarithmische Entwicklung erhalten. Dies würde die numerische Vorhersage für \(\alpha^{-1}\) verändern und muss mit dem CODATA-Wert abgeglichen werden. Vorläufige Schätzungen deuten darauf hin, dass die Verschiebung innerhalb der aktuellen Unsicherheit der topologischen Vorhersage liegt, aber eine genaue Berechnung ist erforderlich und wird in einer zukünftigen Überarbeitung von Anhang~G vorgelegt.
	
	% ============================================================
	\section{Diskussion: Ist YUCT eine Theorie von Allem oder der elektroschwachen Skala?}
	\label{sec:discussion}
	
	Der ursprüngliche YUCT-Rahmen erstrebte es, alle physikalischen Skalen von der Planck-Masse bis zur Elektronenmasse durch eine einzige Koordinationshierarchie zu erklären. Die korrigierte Berechnung der skalaren Masse offenbart eine Spannung in diesem Anspruch: Wenn das Koordinationsnetzwerk auf der Planck-Skala arbeitet, sind die Strahlungskorrekturen enorm und müssen durch feinabgestimmte bloße Massen kompensiert werden.
	
	Szenario B löst diese Spannung, indem es die YUCT-Massenleiter reinterpretiert. Anstatt Massen direkt durch die Tiefe \(L\) zu erzeugen, bestimmt die Leiter die \textbf{Yukawa-Kopplungen}, die wiederum die Massen radiativ auf der Skala \(\LUV \sim 9\) TeV erzeugen. Die Planck-Skala wird dann zu der Skala, auf der das Koordinationsnetzwerk zu \textit{entstehen} beginnt, aber die physikalische Körnung – die Skala, auf der der Formfaktor wichtig wird – ist viel niedriger.
	
	Diese Neuinterpretation bewahrt die algebraische Schönheit von YUCT und beseitigt das Feinabstimmungsproblem. Sie verwandelt YUCT auch von einer spekulativen Theorie der Quantengravitation in ein falsifizierbares Modell der elektroschwachen Physik, das an aktuellen und zukünftigen Collidern getestet werden kann.
	
	% ============================================================
	\section{Schlussfolgerung}
	\label{sec:conclusion}
	
	Wir haben einen Rechenfehler in der skalaren Massenkorrektur innerhalb von YUCT korrigiert und den exakten analytischen Wert \(C_0 = 3\sqrt{\pi}/(8\sqrt{2}) \approx 0,46999\) erhalten. Mit der Planck-Skalen-Identifikation \(\LUV = 3\pi\Mpl\) beträgt die Strahlungskorrektur zur Higgs-Masse \(\delta m_H \approx 1,57\times 10^{18}\) GeV, was eine feinabgestimmte Kompensation mit der bloßen Masse erfordert.
	
	Als Alternative haben wir Szenario B vorgeschlagen, in dem der bloße Lagrange-Operator auf der Planck-Skala skaleninvariant ist und alle Massen radiativ erzeugt werden. Selbstkonsistenz verlangt dann, dass die fundamentale Koordinationsskala \(\LUV \approx 8,96\) TeV beträgt – ein Wert, der am LHC zugänglich ist. Dieses Szenario liefert konkrete, falsifizierbare Vorhersagen für hochenergetische Streuprozesse.
	
	Die Wahl zwischen Szenario A (ästhetisch unbefriedigend, aber konsistent mit dem Planck-Skalen-Ursprung von YUCT) und Szenario B (vorhersagekräftig und testbar, aber eine Neuinterpretation der YUCT-Massenleiter erfordernd) kann nicht allein auf theoretischer Grundlage getroffen werden. Sie muss durch das Experiment entschieden werden. Wir fordern die LHC-Kollaborationen auf, Daten zu Drell-Yan- und Dijet-Ereignissen mit hohen Massen auf Hinweise auf die durch Gl.~(\ref{eq:F}) mit \(\LUV \sim 9\) TeV vorhergesagte Formfaktor-Unterdrückung zu untersuchen. Ein Nullresultat würde Szenario B ausschließen, und obwohl Szenario A logisch möglich bleibt, würde es den Fall für YUCT als Lösung des Hierarchieproblems schwächen.
	
	\section*{Danksagung}
	Der Autor dankt den Teilnehmern des YUCT-Seminars für kritische Diskussionen.
	
	\section*{Datenverfügbarkeit}
	Alle Berechnungen sind eigenständig; der für numerische Auswertungen verwendete Python-Code ist im YPSDC-Repository verfügbar: \url{https://github.com/Alexey-Yakushev-YUCT/YPSDC}.
	
	\begin{thebibliography}{99}
		\bibitem{AppendixL} A.\,V.\,Yakushev, „Anhang L: Fraktale Koordinationsfehlerskalierung,“ in \emph{YUCT}, Zenodo (2026).
		\bibitem{AppendixV} A.\,V.\,Yakushev, „Anhang V: Zeit als Entropie von Koordinationsprozessen,“ in \emph{YUCT}, Zenodo (2026).
		\bibitem{AppendixG} A.\,V.\,Yakushev, „Anhang G: Gravitation als geometrische Spannung des Koordinationsnetzwerks und zyklische Kosmologie in YUCT,“ in \emph{YUCT}, Zenodo (2026).
		\bibitem{AppendixM} A.\,V.\,Yakushev, „Anhang M: YUCT-Kosmologie – Inflation als Koordinationsphasenübergang,“ in \emph{YUCT}, Zenodo (2026).
		\bibitem{AppendixLambda} A.\,V.\,Yakushev, „Anhang \(\Lambda\): Auflösung der Hierarchie- und kosmologischen Konstanten-Probleme in YUCT,“ in \emph{YUCT}, Zenodo (2026).
		\bibitem{AppendixJ} A.\,V.\,Yakushev, „Anhang J: Vollständige Herleitung der Flavour-Parameter des Standardmodells aus topologischen Offsets in YUCT,“ in \emph{YUCT}, Zenodo (2026).
		\bibitem{AppendixFerra} A.\,V.\,Yakushev, „Anhang Ferra1.2: Universelle Koordinationskonstanten für geordnete und ungeordnete Phasen,“ in \emph{YUCT}, Zenodo (2026).
	\end{thebibliography}
	
\end{document}